% Rozwinięcie AF-15 i AF-31, bez teorii widma i rachunku holomorficznego.
\chapter{Algebry Banacha i szereg Neumanna}
\label{ma2:ch:banach-algebras}

Wykorzystamy szeregi w przestrzeniach Banacha oraz
twierdzenia o ciągłych odwzorowaniach dwuliniowych
z rozdziałów \ref{ma2:ch:banach-series}
i \ref{ma1:ch:cauchy-product}.
Podstawy teorii algebr Banacha pochodzą z analizy
funkcjonalnej (rozdziały 15 i 31). Nie potrzebujemy
tutaj teorii widma ani analizy zespolonej.

\section{Algebry unormowane}

\begin{definicja}[Algebra unormowana i algebra Banacha]
\label{ma2:def:normed-algebra}
\label{def:banach-algebra31}
Algebrą unormowaną nad $\K\in\{\R,\C\}$ nazywamy
przestrzeń unormowaną $A$ wyposażoną w łączne
mnożenie dwuliniowe spełniające
\[
 \|ab\|\le\|a\|\|b\|\qquad(a,b\in A).
\]
Jeżeli przestrzeń $A$ jest zupełna, nazywamy ją
algebrą Banacha. Jeżeli istnieje element $1_A$
taki, że $1_Aa=a1_A=a$ dla każdego $a\in A$,
nazywamy go jednostką. W dalszej części rozdziału
zakładamy istnienie jednostki i $\|1_A\|=1$,
chyba że zaznaczono inaczej. Nie zakładamy
przemienności mnożenia.
\end{definicja}
Jest to rozszerzenie definicji zespolonej algebry
Banacha z AF, rozdział 31, podrozdział o elementach
odwracalnych. Mnożenie jest ciągłe na mocy lematu
\ref{ma2:lem:bilinear-bound} ze stałą $K=1$.
Dla każdego $a\in A$ określamy $a^0=1_A$ oraz
$a^{n+1}=a^na$. Dla ustalonego $m\ge0$ równość $a^ma^0=a^m$
jest własnością jednostki. Jeżeli $a^ma^n=a^{m+n}$,
to z łączności
$a^ma^{n+1}=(a^ma^n)a=a^{m+n+1}$.
Zatem $a^ma^n=a^{m+n}$ dla wszystkich $m,n\ge0$.
Ponadto $\|a^0\|=1$, a krok
$\|a^{n+1}\|\le\|a^n\|\|a\|$ dowodzi indukcyjnie
$\|a^n\|\le\|a\|^n$. Stosujemy konwencję
$t^0=1$ również dla $t=0$.

\begin{przyklad}[Podstawowe algebry Banacha]
\label{ex:banach-algebras31}
Niech $E$ będzie niezerową przestrzenią Banacha.
Przestrzeń $\mathcal L(E,E)$ z normą operatorową
i składaniem operatorów jest algebrą Banacha.
Zupełność wynika ze stwierdzenia
\ref{prop:operator-norm12}. Łączność i dwuliniowość
wynikają z własności składania operatorów liniowych.
Dla każdego $x\in E$ mamy
$\|STx\|\le\|S\|\|T\|\|x\|$, skąd
$\|ST\|\le\|S\|\|T\|$. Jednostką jest operator
identycznościowy, którego norma wynosi $1$.

Jeżeli $K$ jest niepustą zwartą przestrzenią metryczną,
to $C(K,\K)$ z normą supremum i mnożeniem punktowym
jest przemienną algebrą Banacha. Zupełność wynika
z twierdzenia \ref{ma2:thm:continuous}, a nierówność
$\|fg\|_\infty\le\|f\|_\infty\|g\|_\infty$ z oszacowania
punktowego. Jednostką jest funkcja stale równa $1$.

Przestrzeń $\ell^1(\N_0,\K)$ ze splotem
\[
 (a*b)_n=\sum_{j=0}^n a_jb_{n-j}
\]
jest przemienną algebrą Banacha. Zupełność po
przeindeksowaniu wynika z twierdzenia
\ref{thm:seq-banach12}. Oszacowanie
$\|a*b\|_1\le\|a\|_1\|b\|_1$
jest dokładnie oszacowaniem z twierdzenia
\ref{ma1:thm:absolute-cauchy-product}.
Dwuliniowość wynika z rozdzielności sum skończonych,
przemienność z zamiany $j$ na $n-j$, a łączność z tego,
że dla każdego $n$ obie współrzędne
$((a*b)*c)_n$ i $(a*(b*c))_n$ są równe
$\sum_{i+j+k=n}a_i b_j c_k$ (indeksy nieujemne).
Wynika to z rozwinięcia skończonych sum
jak w dowodzie wniosku \ref{ma1:cor:finite-cauchy-products}.
Jednostką jest $(1,0,0,\ldots)$ o normie $1$.
\end{przyklad}
W pierwszym przypadku dla $E=\K^d$, $d\ge1$,
otrzymujemy algebrę macierzy z normą operatorową.

\section{Elementy odwracalne}

\begin{definicja}
\label{ma2:def:algebra-invertible}
Element $a\in A$ jest odwracalny, gdy istnieje $b\in A$
takie, że $ab=ba=1_A$. Zbiór takich elementów
oznaczamy przez $A^\times$.
\end{definicja}
Jeżeli $ba=1_A$ i $ac=1_A$, to
$b=b(ac)=(ba)c=c$. W szczególności obustronna
odwrotność jest jednoznaczna; piszemy $a^{-1}$.
Zbiór $A^\times$ jest grupą: jednostka jest odwracalna,
odwrotność elementu odwracalnego ma odwrotność $a$,
a z łączności wynika
$(ab)(b^{-1}a^{-1})=1_A=(b^{-1}a^{-1})(ab)$.
Zatem $(ab)^{-1}=b^{-1}a^{-1}$.

\begin{przyklad}[Odwrotność jednostronna nie wystarcza]
\label{ma2:ex:one-sided-inverse}
W $\ell^2(\N_+)$ określamy przesunięcia
\[
 R(x_1,x_2,\ldots)=(0,x_1,x_2,\ldots),\qquad
 L(x_1,x_2,\ldots)=(x_2,x_3,\ldots).
\]
Są liniowe, $\|Rx\|_2=\|x\|_2$ oraz
$\|Lx\|_2^2=\|x\|_2^2-|x_1|^2\le\|x\|_2^2$,
więc należą do algebry ograniczonych operatorów.
Z definicji $LR=I$, ale $RL(e_1)=0\ne e_1$.
Gdyby $R$ był odwracalny, mnożąc $LR=I$ z prawej
przez $R^{-1}$, otrzymalibyśmy $L=R^{-1}$.
Gdyby $L$ był odwracalny, mnożąc $LR=I$ z lewej
przez $L^{-1}$, otrzymalibyśmy $R=L^{-1}$.
W obu przypadkach wynikałoby $RL=I$, sprzeczność.
\end{przyklad}

\section{Szereg Neumanna}

\begin{twierdzenie}[Szereg Neumanna i oszacowanie reszty]
\label{ma2:thm:neumann}
\label{lem:neumann31}
\label{lem:neumann15}
Niech $A$ będzie algebrą Banacha z jednostką
o normie $1$. Jeśli $u\in A$ i $q=\|u\|<1$, to
\[
 (1_A-u)^{-1}=\sum_{n=0}^\infty u^n,\qquad
 \|(1_A-u)^{-1}\|\le\frac1{1-q}.
\]
Dla $N\ge0$ mamy ponadto
\[
 \left\|(1_A-u)^{-1}-\sum_{n=0}^N u^n\right\|
 \le\frac{q^{N+1}}{1-q}.
\]
\end{twierdzenie}
\begin{proof}
Twierdzenie w przypadku zespolonej algebry Banacha
i jego wersja operatorowa zostały udowodnione
w AF (odpowiednio rozdziały 31 i 15).
Podajemy krótki argument wspólny dla obu ciał,
z oszacowaniem reszty. Z $\|u^n\|\le q^n$ oraz
twierdzenia \ref{ma2:thm:absolute-series}
(po przeindeksowaniu) szereg ma sumę $v\in A$.
Dla $S_N=\sum_{n=0}^N u^n$ rozdzielność daje
$(1_A-u)S_N=S_N(1_A-u)=1_A-u^{N+1}$.
Z $\|u^{N+1}\|\le q^{N+1}\to0$ i ciągłości
mnożenia otrzymujemy $(1_A-u)v=v(1_A-u)=1_A$.
Z nierówności trójkąta $\|S_N\|\le\sum_{n=0}^N q^n$
i ciągłości normy otrzymujemy $\|v\|\le1/(1-q)$.
Oszacowanie reszty wynika z twierdzenia
\ref{ma2:thm:absolute-series}, ponieważ
$\sum_{n>N}q^n=q^{N+1}/(1-q)$ na mocy wzoru
na szereg geometryczny (AM I, rozdział 09).
Dla $q=0$ mamy $u=0$, a wszystkie reszty są zerowe.
\end{proof}

\begin{wniosek}[Słabsze kryterium potęgowe]
\label{ma2:cor:neumann-power}
Dla $u\in A$ równoważne są następujące warunki:
\begin{enumerate}[label=(\roman*)]
\item szereg $\sum_{n\ge0}u^n$ jest zbieżny;
\item szereg ten jest bezwzględnie zbieżny;
\item istnieje $m\ge1$ takie, że $\|u^m\|<1$.
\end{enumerate}
Jeżeli zachodzą, suma szeregu jest równa
$(1_A-u)^{-1}$.
\end{wniosek}
\begin{proof}
(i) implikuje (iii), ponieważ wyrazy zbieżnego
szeregu dążą do zera (twierdzenie
\ref{ma2:thm:series-cauchy}).
Niech zachodzi (iii), $q_m=\|u^m\|<1$.
Każdy $n\ge0$ ma jednoznaczny rozkład
$n=km+r$ z $k\ge0$ i $0\le r<m$.
Z podmultiplikatywności
$\|u^{km+r}\|\le q_m^k\|u^r\|$.
Dla każdej skończonej sumy częściowej sumę norm
ograniczamy przez
$(\sum_{r=0}^{m-1}\|u^r\|)\sum_{k\ge0}q_m^k$.
Jest to skończona liczba, więc zachodzi (ii).
(ii) implikuje (i) z twierdzenia
\ref{ma2:thm:absolute-series}.
Ponieważ $u^{N+1}\to0$, tożsamość teleskopowa
z dowodu twierdzenia \ref{ma2:thm:neumann}
daje obustronną odwrotność.
\end{proof}

\begin{przyklad}
\label{ma2:ex:neumann-nilpotent}
Dla macierzy $U=\begin{pmatrix}0&2\\0&0\end{pmatrix}$
mamy $U^2=0$ oraz $\|U\|=2$ w normie operatorowej
indukowanej przez normę euklidesową (dla
$(x,y)$ obraz ma normę $2|y|$).
Mimo $\|U\|>1$ szereg Neumanna kończy się po
dwóch składnikach i $(I-U)^{-1}=I+U$.
Natomiast dla $u=2$ w $\R$ element $1-u=-1$
jest odwracalny, ale szereg $\sum 2^n$ jest rozbieżny.
Sama odwracalność $1-u$ nie wystarcza do zbieżności
szeregu Neumanna.
\end{przyklad}

\section{Zaburzenia elementów odwracalnych}

\begin{twierdzenie}[Oszacowania odwrotności]
\label{ma2:thm:inverse-perturbation}
Niech $a\in A^\times$ i $h\in A$.
Jeżeli $t=\|a^{-1}h\|<1$, to $a+h\in A^\times$,
\[
 (a+h)^{-1}=\sum_{n=0}^\infty(-a^{-1}h)^n a^{-1},
 \qquad
 \|(a+h)^{-1}\|\le\frac{\|a^{-1}\|}{1-t}.
\]
Ponadto
\[
 \|(a+h)^{-1}-a^{-1}\|\le
 \frac{t\|a^{-1}\|}{1-t}.
\]
Reszta po wyrazie rzędu $N$ jest nie większa
niż $t^{N+1}\|a^{-1}\|/(1-t)$.
\end{twierdzenie}
\begin{proof}
Mamy $a+h=a(1_A+a^{-1}h)$.
Z twierdzenia \ref{ma2:thm:neumann} zastosowanego
do $u=-a^{-1}h$ drugi czynnik jest odwracalny.
Reguła odwrotności iloczynu daje rozwinięcie
z zachowaniem kolejności czynników. Ciągłość
mnożenia z prawej strony przez $a^{-1}$
pozwala pomnożyć granicę szeregu. Wszystkie
oszacowania wynikają z twierdzenia
\ref{ma2:thm:neumann} i nierówności
$\|va^{-1}\|\le\|v\|\|a^{-1}\|$;
w oszacowaniu różnicy bierzemy $N=0$.
\end{proof}
W szczególności, ponieważ $a^{-1}\ne0$
(gdyż $aa^{-1}=1_A$ i $\|1_A\|=1$), kula otwarta
o środku $a$ i promieniu $1/\|a^{-1}\|$
mieści się w $A^\times$,
gdyż $\|a^{-1}h\|\le\|a^{-1}\|\|h\|$.
Zatem $A^\times$ jest otwarty. Jeżeli
$\|h\|\le1/(2\|a^{-1}\|)$, to
$\|(a+h)^{-1}-a^{-1}\|\le2\|a^{-1}\|^2\|h\|$.
Otrzymujemy ciągłość odwracania, zgodnie
z wynikiem AF-31 (lemat o szeregu Neumanna).

\begin{wniosek}[Pierwsze przybliżenie odwrotności]
\label{ma2:cor:inverse-linear}
Przy $t=\|a^{-1}h\|<1$ zachodzi
\[
 \|(a+h)^{-1}-a^{-1}+a^{-1}ha^{-1}\|
 \le\frac{t^2\|a^{-1}\|}{1-t}.
\]
Jeśli $\|h\|\le 1/(2\|a^{-1}\|)$, prawa strona
jest nie większa niż $2\|a^{-1}\|^3\|h\|^2$.
\end{wniosek}
\begin{proof}
Bierzemy $N=1$ w poprzednim twierdzeniu.
Drugie oszacowanie wynika z $t\le1/2$ oraz
$t\le\|a^{-1}\|\|h\|$.
\end{proof}
Jest to przygotowanie do rachunku różniczkowego:
liniowa część przyrostu odwrotności to
$h\mapsto-a^{-1}ha^{-1}$. Pojęcie pochodnej
Frécheta i pełna teoria różniczkowania
odwracania znajdują się w AF, rozdział 15.

\section{Szeregi potęgowe w algebrze Banacha}

\begin{twierdzenie}[Podstawianie elementu do szeregu]
\label{ma2:thm:algebra-power-series}
Niech $(c_n)_{n\ge0}$ będzie ciągiem skalarów.
Jeżeli dla $r\ge0$ zachodzi $\sum_{n\ge0}|c_n|r^n<\infty$,
to dla każdego $a\in A$ z $\|a\|\le r$ istnieje
$f(a)=\sum_{n\ge0}c_na^n$ oraz
\[
 \left\|f(a)-\sum_{n=0}^N c_na^n\right\|
 \le\sum_{n>N}|c_n|r^n.
\]
Jeżeli $(d_n)_{n\ge0}$ jest drugim ciągiem skalarów
i $\sum_{n\ge0}|d_n|r^n<\infty$, to definiując
$g(a)=\sum_{n\ge0}d_na^n$, dla $\|a\|\le r$
otrzymujemy bezwzględnie zbieżny szereg
\[
 f(a)g(a)=\sum_{n=0}^\infty
     \left(\sum_{j=0}^n c_jd_{n-j}\right)a^n.
\]
\end{twierdzenie}
\begin{proof}
Z $\|c_na^n\|\le |c_n|r^n$ i twierdzenia
\ref{ma2:thm:absolute-series} otrzymujemy
istnienie sumy oraz oszacowanie reszty.
Do obu szeregów w $A$ stosujemy wniosek
\ref{ma2:cor:bilinear-absolute} dla mnożenia
algebry. Skończona suma współczynników
iloczynu Cauchy'ego ma podaną postać, ponieważ
$(c_ja^j)(d_{n-j}a^{n-j})=c_jd_{n-j}a^n$.
Nie potrzebujemy przemienności całej algebry.
\end{proof}

\begin{wniosek}[Funkcja wykładnicza]
\label{ma2:cor:algebra-exp}
Dla każdego $a\in A$ istnieje
$\exp a=\sum_{n\ge0}a^n/n!$ i $\|\exp a\|\le e^{\|a\|}$.
Jeżeli $a,b$ komutują, to
$\exp(a+b)=(\exp a)(\exp b)$.
W szczególności $\exp a$ jest odwracalny i
$(\exp a)^{-1}=\exp(-a)$.
\end{wniosek}
\begin{proof}
Szereg skalarny $\sum r^n/n!$ ma sumę $e^r$
z definicji funkcji wykładniczej w AM I,
rozdział 18. Twierdzenie
\ref{ma2:thm:algebra-power-series} zapewnia
istnienie sumy; oszacowanie normy wynika
z nierówności trójkąta i przejścia do granicy.
Stosujemy wniosek \ref{ma2:cor:bilinear-absolute}
do szeregów o wyrazach $a^n/n!$ i $b^n/n!$.
Przy $ab=ba$ wzór dwumianowy daje
\[
 \sum_{j=0}^n\frac{a^jb^{n-j}}{j!(n-j)!}
 =\frac{(a+b)^n}{n!}.
\]
Wzór dwumianowy w tym przypadku wynika przez
indukcję: mnożymy rozwinięcie dla $n$ przez
$a+b$, przesuwamy $a$ przez potęgi $b$ dzięki
przemienności pary i łączymy współczynniki
regułą Pascala. Rząd $0$ daje jednostkę.
Stąd prawo dodawania argumentów dla
komutujących elementów. Para $a,-a$ komutuje,
a $\exp0=1_A$, więc otrzymujemy obustronną odwrotność.
\end{proof}

\section{Przykład obliczeniowy}

\begin{przyklad}[Rozwiązanie równania operatorowego]
\label{ma2:ex:neumann-equation}
Niech $E$ będzie niezerową przestrzenią Banacha,
$T$ ograniczonym operatorem liniowym na $E$
i $0\le\|T\|\le q<1$.
Dla $b\in E$ równanie $x-Tx=b$ ma dokładnie jedno
rozwiązanie
$x=\sum_{n\ge0}T^nb$.
Istotnie, w algebrze operatorów $I-T$ jest odwracalny
z twierdzenia \ref{ma2:thm:neumann}.
Ewaluacja na $b$ jest ciągła, bo
$\|Sb\|\le\|S\|\|b\|$, więc rozwinięcie odwrotności
daje wskazaną sumę. Dla
$x_N=\sum_{n=0}^N T^nb$ mamy
\[
 \|x-x_N\|\le\frac{q^{N+1}}{1-q}\|b\|,
 \qquad b-(I-T)x_N=T^{N+1}b.
\]
Drugą równość daje tożsamość teleskopowa.
Dla dowolnego przybliżenia $y\in E$ i jego reszty
$r=b-(I-T)y$ zachodzi ponadto
\[
 x-y=(I-T)^{-1}r,\qquad
 \|x-y\|\le\frac{\|r\|}{1-q}.
\]
Pierwsza równość wynika z $(I-T)(x-y)=r$,
a druga z oszacowania normy odwrotności.
Jest to oszacowanie błędu na podstawie obliczonej
reszty, niezależne od sposobu uzyskania $y$.
W szczególności dla $q=1/2$ i $N=9$
błąd normy nie przekracza $2^{-9}\|b\|$.
Dla $E=\{0\}$ równanie ma jedyne rozwiązanie $0$
i wszystkie oszacowania są trywialnie spełnione.
\end{przyklad}

W kolejnym rozdziale zajmiemy się zbieżnością
punktową i jednostajną ciągów funkcji.
